% =====================================================================
%  exposicion_clients.tex
%  Agregador resiliente de APIs - Clientes de dependencias externas
%  Documento de apoyo para la exposicion (parte de Emerson)
%  Compilar con:  pdflatex exposicion_clients.tex  (dos veces)
% =====================================================================
\documentclass[11pt,a4paper]{article}

% ---------- Codificacion e idioma ----------
\usepackage[utf8]{inputenc}
\usepackage[T1]{fontenc}
\usepackage[spanish,es-noquoting,es-nodecimaldot]{babel}
\usepackage{lmodern}
\usepackage{textcomp}
\usepackage{amsmath,amssymb}

% ---------- Geometria y tipografia ----------
\usepackage[margin=2.4cm]{geometry}
\usepackage{microtype}
\usepackage{parskip}
\emergencystretch=3em

% ---------- Tablas, figuras, graficos ----------
\usepackage{booktabs}
\usepackage{array}
\usepackage{graphicx}
\usepackage{float}
\usepackage{tikz}
\usetikzlibrary{arrows.meta,positioning,shapes.geometric,babel}
\usepackage{xcolor}

% ---------- Codigo fuente ----------
\usepackage{listings}
\usepackage{upquote}

% ---------- Hipervinculos ----------
\usepackage[hidelinks]{hyperref}

% ---------- Colores ----------
\definecolor{azul}{HTML}{1F4E79}
\definecolor{verde}{HTML}{2E7D32}
\definecolor{gris}{HTML}{F2F2F2}
\definecolor{grisborde}{HTML}{BDBDBD}
\definecolor{naranja}{HTML}{B45309}

\hypersetup{
  colorlinks=true,
  linkcolor=azul,
  urlcolor=azul,
  citecolor=azul,
  pdftitle={Agregador resiliente de APIs - Clientes externos},
  pdfauthor={Emerson}
}

% ---------- Estilo de listings (Python) ----------
\lstdefinestyle{py}{
  language=Python,
  basicstyle=\ttfamily\footnotesize,
  keywordstyle=\color{azul}\bfseries,
  commentstyle=\color{verde}\itshape,
  stringstyle=\color{naranja},
  numberstyle=\tiny\color{gray},
  numbers=left,
  numbersep=8pt,
  stepnumber=1,
  frame=single,
  rulecolor=\color{grisborde},
  backgroundcolor=\color{gris},
  showstringspaces=false,
  breaklines=true,
  breakatwhitespace=false,
  tabsize=4,
  upquote=true,
  extendedchars=true,
  inputencoding=utf8,
  columns=fullflexible,
  keepspaces=true,
  literate=%
    {á}{{\'a}}1 {é}{{\'e}}1 {í}{{\'i}}1 {ó}{{\'o}}1 {ú}{{\'u}}1
    {Á}{{\'A}}1 {É}{{\'E}}1 {Í}{{\'I}}1 {Ó}{{\'O}}1 {Ú}{{\'U}}1
    {ñ}{{\~n}}1 {Ñ}{{\~N}}1 {ü}{{\"u}}1
    {°}{{$^{\circ}$}}1 {¿}{{?`}}1 {¡}{{!`}}1 {–}{{-}}1 {—}{{---}}1
}
\lstset{style=py}

% ---------- Macros de comodidad ----------
\newcommand{\arch}[1]{\texttt{#1}}
\newcommand{\quot}{\textquotedbl}
\newcommand{\req}[1]{\textbf{\textcolor{azul}{#1}}}
\newcommand{\ok}{\textcolor{verde}{\textbf{\checkmark}}}
\newcommand{\no}{\textcolor{red}{\textbf{$\times$}}}

% =====================================================================
\begin{document}

% ---------------------------- PORTADA --------------------------------
\begin{titlepage}
  \centering
  \vspace*{2.2cm}
  {\scshape\Large Sistemas Distribuidos \par}
  \vspace{0.4cm}
  {\large Taller: Agregador resiliente de APIs \par}
  \vspace{2.5cm}
  \rule{\linewidth}{0.5pt}\par
  \vspace{0.6cm}
  {\Huge\bfseries Clientes de dependencias externas \par}
  \vspace{0.4cm}
  {\Large \arch{app/clients} \; y \; \arch{tests/clients} \par}
  \vspace{0.6cm}
  \rule{\linewidth}{0.5pt}\par
  \vspace{2.2cm}

  {\large \textbf{Autor (esta parte):} Emerson \par}
  \vspace{0.2cm}
  {\large \textbf{Rol:} consumo resiliente de dependencias \par}
  \vspace{1.6cm}

  \begin{tabular}{rl}
    \toprule
    \textbf{Paquete} & \arch{app/clients/} \\
    \textbf{Puebas}  & \arch{tests/clients/} \\
    \textbf{API externa elegida} & USGS (terremotos) \\
    \textbf{Servicio inestable} & \arch{servicio\_inestable.py} (puerto 9000) \\
    \textbf{Estado} & 20 pruebas en verde \\
    \bottomrule
  \end{tabular}

  \vfill
  {\large \today \par}
\end{titlepage}

\tableofcontents
\newpage

% =====================================================================
\section{Introduccion y contexto}
% =====================================================================

El taller pide construir una \textbf{API agregadora resiliente}: un servicio
propio (FastAPI) que \emph{agrega} dos dependencias que no controlamos:

\begin{enumerate}
  \item un \textbf{servicio local inestable} (sensores) que falla de forma
        aleatoria (503), responde lentamente, corta conexiones y aplica
        limites de tasa (429 con \arch{Retry-After});
  \item una \textbf{API publica real} por Internet (en este caso, los
        terremotos de \textbf{USGS}).
\end{enumerate}

Mi responsabilidad es unicamente la capa de \textbf{clientes}: la pieza que
sabe \emph{hablar} con esas dos dependencias y \emph{sobrevivir} cuando se
portan mal. El resto del proyecto lo cubren mis companeras:

\begin{itemize}
  \item \textbf{Nataly}: \arch{app/config.py}, \arch{app/models.py},
        \arch{app/cache.py}.
  \item \textbf{Daniela}: \arch{app/routers/*}.
  \item \textbf{Grupal}: \arch{main.py}, \arch{README.md},
        \arch{requirements.txt}.
\end{itemize}

\subsection{Que se evalua (rubrica)}
\begin{table}[H]
\centering
\begin{tabular}{lr l}
\toprule
\textbf{Criterio} & \textbf{Puntos} & \textbf{Donde cae} \\
\midrule
Consumo de dependencias          & 15 & \arch{app/clients/} \\
API propia (endpoints)           & 20 & routers (Daniela) \\
Resiliencia                      & 30 & \arch{app/clients/} \\
Cache + latencia                 & 10 & \arch{app/clients/} \\
Idempotencia                     & 10 & routers (Daniela) \\
Analisis / hallazgos             & 15 & informe (grupo) \\
\midrule
Bonus                            & +10 & extras bien justificados \\
\bottomrule
\end{tabular}
\caption{Reparto de la rubrica y de quien lo cubre.}
\end{table}

% =====================================================================
\section{Arquitectura general}
% =====================================================================

La figura~\ref{fig:arq} ubica mi parte dentro del sistema completo. Las
flechas discontinuas son dependencias que \emph{yo no controlo}: por eso el
cliente debe ser defensivo.

\begin{figure}[H]
\centering
\begin{tikzpicture}[
  font=\small,
  caja/.style={rectangle, rounded corners, draw, thick, align=center,
               minimum height=1.1cm, minimum width=2.9cm, fill=white},
  mio/.style={caja, fill=azul!12, draw=azul, very thick},
  ext/.style={caja, fill=orange!12, draw=naranja, dashed, very thick},
  >={Stealth[length=2.5mm]}
]
  \node[mio] (clients) {%
    \textbf{app/clients}\\[-2pt]
    \footnotesize Emerson\\[-2pt]
    \scriptsize (resiliencia)};
  \node[caja, below=1.5cm of clients, xshift=-3.4cm] (main)
    {main.py\\\scriptsize FastAPI};
  \node[caja, below=1.5cm of clients, xshift=3.4cm] (routers)
    {routers / v1\\\scriptsize Daniela};

  \node[ext, above left=1.4cm and 1.0cm of clients] (inestable)
    {servicio\_inestable\\\scriptsize :9000 (falla)};
  \node[ext, above right=1.4cm and 1.0cm of clients] (usgs)
    {USGS (HTTPS)\\\scriptsize Internet};

  \node[caja, left=1.4cm of clients, fill=verde!10, draw=verde] (cache)
    {cache.py\\\scriptsize Nataly};
  \node[caja, right=1.4cm of clients, fill=verde!10, draw=verde] (models)
    {models.py\\\scriptsize Nataly};

  \draw[->, thick] (main) -- (clients);
  \draw[->, thick] (routers) -- (clients);
  \draw[<->, dashed, naranja] (clients) -- (inestable);
  \draw[<->, dashed, naranja] (clients) -- (usgs);
  \draw[->, dashed, verde] (clients) -- (cache);
  \draw[->, dashed, verde] (clients) -- (models);
\end{tikzpicture}
\caption{Arquitectura general. En azul, mi parte; en naranja punteado, las
dependencias que no controlo; en verde, los modulos de Nataly.}
\label{fig:arq}
\end{figure}

% =====================================================================
\section{Mi alcance y el contrato}
% =====================================================================

Solo toco estos archivos (nada de mis companeras):

\begin{table}[H]
\centering
\small
\begin{tabular}{p{4.4cm} p{9.4cm}}
\toprule
\textbf{Archivo} & \textbf{Responsabilidad} \\
\midrule
\arch{\_\_init\_\_.py} &
  Jerarquia de excepciones del contrato. Se define aqui (y no en los
  submodulos) para evitar importaciones circulares. \\
\arch{unstable\_service.py} &
  Cliente resiliente del servicio de sensores: timeouts, reintentos
  selectivos, backoff, \arch{Retry-After}, bloqueo por IP, paginacion,
  degradacion parcial. \\
\arch{external\_api.py} &
  Cliente de USGS: dos llamadas encadenadas (feed $\rightarrow$ detalle),
  validacion Pydantic, cache con TTL y medicion de latencia. \\
\arch{tests/clients/} &
  Pruebas automatizadas (pytest + respx) y evidencia real. \\
\bottomrule
\end{tabular}
\caption{Archivos de mi parte (dentro de \arch{app/clients/}).}
\end{table}

\subsection{Decision de diseno: API por clases}
El esqueleto original define los clientes como \textbf{clases} y
\arch{main.py} / el router de Daniela ya las instancian asi
(\arch{ExternalEarthquakeClient}, \arch{UnstableServiceClient}).
Por eso \textbf{mantengo las clases} y conservo sus firmas, y ademas expongo
funciones equivalentes del contrato (\arch{get\_con\_reintentos},
\arch{obtener\_ultimo\_sismo}). Cambiar a un estilo solo-funciones habria
roto el codigo de mis companeras.

% =====================================================================
\section{Jerarquia de excepciones (\arch{\_\_init\_\_.py})}
% =====================================================================

El router traduce las excepciones a codigos HTTP. Para que eso funcione sin
tocar su codigo, las excepciones nuevas \textbf{heredan} de las que ya se
usaban (\arch{UnstableServiceError} y \arch{ExternalApiError}) y ademas de
\arch{TimeoutError} cuando corresponde (para que el 504 se dispare solo).

\begin{figure}[H]
\centering
\begin{tikzpicture}[
  font=\small,
  exc/.style={rectangle, draw, thick, rounded corners, align=center,
              minimum height=0.95cm, fill=white, inner sep=4pt},
  >={Stealth[length=2.5mm]}
]
  \node[exc, fill=red!12] (base) {\arch{Exception}};
  \node[exc, fill=azul!15, below=of base] (dep)
    {\arch{ErrorDependencia}\\\scriptsize (router $\to$ 502)};
  \node[exc, below left=1.2cm and 0.4cm of dep] (uns)
    {\arch{UnstableServiceError}\\\scriptsize compatibilidad};
  \node[exc, below right=1.2cm and 0.4cm of dep] (ext)
    {\arch{ExternalApiError}\\\scriptsize compatibilidad};
  \node[exc, below=1.1cm of uns, fill=orange!15] (to)
    {\arch{TimeoutDependencia}\\{\scriptsize (y \arch{TimeoutError})}\\\scriptsize router $\to$ 504};
  \node[exc, below left=1.1cm and -0.2cm of to] (noe)
    {\arch{NoEncontrado}\\\scriptsize $\to$ 404};
  \node[exc, below right=1.1cm and -0.2cm of to] (inv)
    {\arch{RespuestaInvalida}\\\scriptsize $\to$ 502};

  \draw[->, thick] (base) -- (dep);
  \draw[->, thick] (dep) -- (uns);
  \draw[->, thick] (dep) -- (ext);
  \draw[->, thick] (uns) -- (to);
  \draw[->, thick] (uns) -- (noe);
  \draw[->, thick] (uns) -- (inv);
\end{tikzpicture}
\caption{Jerarquia de excepciones. Las nuevas heredan de las de
compatibilidad, asi el router sigue funcionando.}
\label{fig:exc}
\end{figure}

\begin{lstlisting}[caption={app/clients/\_\_init\_\_.py (fragmento)}]
class ErrorDependencia(Exception):
    """Fallo de una dependencia que no se puede resolver (el router -> 502)."""

class UnstableServiceError(ErrorDependencia):
    """Error del servicio inestable. Compatibilidad con main.py/router (502)."""

class ExternalApiError(ErrorDependencia):
    """Error de la API externa. Compatibilidad con main.py/router (502)."""

class TimeoutDependencia(UnstableServiceError, TimeoutError):
    """La dependencia no respondio en el tiempo permitido (router -> 504)."""

class NoEncontrado(UnstableServiceError):
    """El recurso no existe en la dependencia (404 de origen -> 404)."""

class RespuestaInvalida(UnstableServiceError):
    """La respuesta no cumple el modelo Pydantic esperado (router -> 502)."""
\end{lstlisting}

\textbf{Idea clave para explicar:} \arch{TimeoutDependencia} hereda de
\arch{TimeoutError}; el router atrapa \arch{asyncio.TimeoutError} para el 504,
asi que un timeout real del cliente cae \emph{naturalmente} en el 504.

% =====================================================================
\section{Cliente resiliente (\arch{unstable\_service.py})}
% =====================================================================

Este archivo concentra el \textbf{30\% de la rubrica}. Cada mecanismo es un
requisito independiente.

\begin{table}[H]
\centering
\begin{tabular}{c p{6.2cm} p{5.6cm}}
\toprule
\textbf{Req.} & \textbf{Mecanismo} & \textbf{Evidencia en codigo} \\
\midrule
C2 & Timeout explicito (conexion/lectura separadas) & \arch{\_tiempo\_timeout} \\
C3 & Reintentos selectivos (solo 503/429/timeout) & \arch{get\_con\_reintentos} \\
C4 & Backoff exponencial + \emph{full jitter} & \arch{\_espera\_backoff} \\
C5 & \arch{Retry-After} + bloqueo compartido por IP & \arch{\_marcar\_bloqueo}, \arch{\_no\_antes\_de} \\
C6 & Limite de concurrencia & \arch{\_semaforo = Semaphore(3)} \\
C7 & Paginacion que sigue \arch{siguiente} & \arch{get\_all\_sensors} \\
--- & Degradacion parcial & \arch{obtener\_lecturas} \\
--- & Logging de cada intento & \arch{\_registrar} \\
\bottomrule
\end{tabular}
\caption{Requisitos cubiertos por el cliente resiliente.}
\end{table}

\subsection{C2 - Timeout explicito}
La lectura se acota a \textbf{2\,s} deliberadamente \emph{por debajo} de la
lentitud simulada (4\,s): asi una respuesta lenta cuenta como \emph{timeout},
que si es reintentable.

\begin{lstlisting}[caption={Timeout separado conexion/lectura}]
TIMEOUT_CONEXION = 2.0
TIMEOUT_LECTURA = 2.0   # menor que los 4 s de lentitud simulada

def _tiempo_timeout(timeout_seconds: float) -> httpx.Timeout:
    return httpx.Timeout(
        connect=TIMEOUT_CONEXION,
        read=min(timeout_seconds, TIMEOUT_LECTURA),
        write=timeout_seconds,
        pool=timeout_seconds,
    )
\end{lstlisting}

\subsection{C3 - Reintentos selectivos}
Solo se reintentan \textbf{503}, \textbf{429} y \textbf{timeouts}. Los
errores legitimas de cliente (400, 404) se devuelven tal cual para que la
capa superior los traduzca (404 $\to$ 404).

\begin{lstlisting}[caption={Bucle de reintentos selectivos}]
for intento in range(1, max_intentos + 1):
    await _esperar_turno()           # respeta el bloqueo compartido tras 429
    try:
        async with _semaforo:        # maximo 3 peticiones simultaneas
            respuesta = await client.get(url, params=params, timeout=timeout)
    except httpx.TimeoutException as exc:
        if intento >= max_intentos:
            raise TimeoutDependencia(...) from exc
        espera, origen = _espera_backoff(intento, backoff_base)
        _registrar(intento, max_intentos, "timeout", espera, origen)
        await _dormir(espera)
        continue

    if respuesta.status_code == 429: ...
    if respuesta.status_code == 503: ...

    return respuesta   # 400, 404 y demas: NO se reintentan
\end{lstlisting}

\subsection{C4 - Backoff exponencial con \emph{full jitter}}
La espera es aleatoria entre 0 y el tope exponencial. El jitter evita que
muchas tareas reintenten \emph{a la vez} (efecto manada).

\begin{lstlisting}[caption={Calculo de la espera (full jitter)}]
def _espera_backoff(intento: int, base: float) -> tuple[float, str]:
    tope = min(BACKOFF_TOPE, base * (2 ** (intento - 1)))
    return random.uniform(0.0, tope), "backoff"
\end{lstlisting}

\subsection{C5 - \arch{Retry-After} y bloqueo compartido por IP}
Si el servicio responde 429 con \arch{Retry-After}, esa espera
\textbf{manda} sobre el backoff. Ademas se marca un instante global
\arch{\_no\_antes\_de}: \emph{todas} las tareas lo respetan (el servicio
bloquea por IP, no por peticion).

\begin{lstlisting}[caption={Bloqueo global por IP}]
_lock_espera = asyncio.Lock()
_no_antes_de = 0.0

async def _marcar_bloqueo(segundos: float) -> None:
    global _no_antes_de
    async with _lock_espera:
        _no_antes_de = max(_no_antes_de, time.monotonic() + segundos)

async def _esperar_turno() -> None:
    async with _lock_espera:
        restante = _no_antes_de - time.monotonic()
    if restante > 0:
        await _dormir(restante)
\end{lstlisting}

\subsection{C6 - Limite de concurrencia}
Un semaforo evita saturar al servicio (y provocar mas 429).

\begin{lstlisting}[caption={Semaforo de concurrencia}]
_semaforo = asyncio.Semaphore(3)   # maximo 3 peticiones a la vez
\end{lstlisting}

\subsection{C7 - Paginacion con reintento por pagina}
Se sigue \arch{siguiente} hasta \arch{null}. El reintento ocurre
\emph{dentro de cada pagina}: que falle una no obliga a repetir todo.

\begin{lstlisting}[caption={Paginacion completa}]
async def get_all_sensors(self) -> list[Sensor]:
    sensores: list[Sensor] = []
    cursor: str | None = None
    while True:
        pagina = await self._pedir_pagina(cursor)   # reintenta por pagina
        sensores.extend(
            Sensor(id=crudo.id, ubicacion=crudo.ubicacion)
            for crudo in pagina.sensores
        )
        if pagina.siguiente is None:
            return sensores
        cursor = pagina.siguiente
\end{lstlisting}

\subsection{Degradacion parcial y logging}
\arch{obtener\_lecturas} pide varios sensores en paralelo y \textbf{no deja
que uno malo tumbe al resto}: devuelve la dupla (lecturas ok, fallidos con
motivo). Cada intento se registra con \arch{logging} (causa, espera y
origen), lo que produce la evidencia de reintentos.

\begin{lstlisting}[caption={Degradacion parcial}]
resultados = await asyncio.gather(
    *(self.obtener_lectura(sid) for sid in sensor_ids),
    return_exceptions=True,
)
for sensor_id, resultado in zip(sensor_ids, resultados):
    if isinstance(resultado, Lectura):
        lecturas.append(resultado)
    else:
        fallidos.append({"sensor_id": sensor_id, "motivo": str(resultado)})
\end{lstlisting}

% =====================================================================
\section{Cliente de la API externa (\arch{external\_api.py})}
% =====================================================================

USGS se consume con \textbf{dos llamadas encadenadas}:
\begin{enumerate}
  \item el \emph{feed} \arch{all\_hour} (\arch{UsgsFeed}) para hallar el
        evento mas reciente;
  \item el \emph{detalle} de ese evento (\arch{UsgsFeature}).
\end{enumerate}

\begin{figure}[H]
\centering
\begin{tikzpicture}[
  font=\small, node distance=1.1cm,
  caja/.style={rectangle, rounded corners, draw, thick, align=center,
               minimum height=1cm, minimum width=2.6cm, fill=white},
  >={Stealth[length=2.5mm]}
]
  \node[caja, fill=azul!12] (cli) {Cliente\\\scriptsize Emerson};
  \node[caja, right=of cli] (feed) {feed all\_hour\\\scriptsize USGS};
  \node[caja, right=of feed] (det) {detalle\\\scriptsize del evento};
  \node[caja, right=of det, fill=verde!12] (cache) {cache TTL\\\scriptsize 300 s};
  \draw[->, thick] (cli) -- node[above]{\scriptsize 1} (feed);
  \draw[->, thick] (feed) -- node[above]{\scriptsize 2} (det);
  \draw[->, thick] (det) -- (cache);
\end{tikzpicture}
\caption{Consulta encadenada feed $\rightarrow$ detalle, con cache TTL.}
\label{fig:usgs}
\end{figure}

\subsection{Validacion Pydantic y el hallazgo de \arch{extra=\quot forbid\quot}}
Los modelos de \arch{models.py} (de Nataly) declaran \arch{extra=\quot forbid\quot}.
Pero el feed \textbf{real} de USGS trae campos adicionales
(\arch{bbox}, \arch{tz}, \arch{sig}, \arch{magType}, \arch{felt}\ldots).
Resultado: contra la API real, la validacion fallaba y el endpoint habria
devuelto \textbf{502}.

\textbf{Solucion dentro de mi alcance} (sin tocar \arch{models.py}): antes de
validar, se \emph{recortan} las claves desconocidas usando los propios campos
del modelo. Sigo validando con los modelos de Nataly, pero tolero los extras.

\begin{lstlisting}[caption={Tolerancia a campos extra de USGS}]
def _recortar(modelo: type[BaseModel], datos: dict) -> dict:
    """Deja solo los campos que declara el modelo."""
    return {c: v for c, v in datos.items() if c in modelo.model_fields}

def _feature_normalizado(datos: dict) -> dict:
    limpio = _recortar(UsgsFeature, datos)
    limpio["properties"] = _recortar(UsgsProperties, datos.get("properties") or {})
    geo = datos.get("geometry")
    limpio["geometry"] = _recortar(UsgsGeometry, geo) if geo else None
    return limpio
\end{lstlisting}

\subsection{Cache TTL y resultado enriquecido}
\arch{get\_latest\_event} devuelve un \arch{ExternalResult} con el sismo, si
fue \emph{hit} de cache y la latencia medida. \arch{usar\_cache=False} fuerza
la llamada real (necesario para medir latencias).

\begin{lstlisting}[caption={Consulta con cache TTL}]
@dataclass(frozen=True, slots=True)
class ExternalResult:
    feature: UsgsFeature | None
    cache_hit: bool
    latency_ms: float

async def get_latest_event(self, usar_cache: bool = True) -> ExternalResult:
    inicio = time.perf_counter()
    if usar_cache:
        guardado = await self._cache.get(self._feed_url)
        if guardado is not None:
            return ExternalResult(guardado, True,
                                  (time.perf_counter() - inicio) * 1000)
    sismo = await _encadenar_usgs(self._client, self._feed_url, self._timeout)
    if sismo is not None:
        await self._cache.set(self._feed_url, sismo)
    return ExternalResult(sismo, False, (time.perf_counter() - inicio) * 1000)
\end{lstlisting}

\subsection{C8 - Medicion de latencia (media y p95)}
Se calienta la cache y luego se mide: \textbf{con cache} (hits) y
\textbf{sin cache} (llamadas reales). El percentil es por
\emph{rango mas cercano} (\emph{nearest-rank}), simple de explicar.

\begin{lstlisting}[caption={Percentil y resumen de latencia}]
def _percentil(valores: list[float], porcentaje: float) -> float:
    if not valores:
        return 0.0
    ordenados = sorted(valores)
    indice = ceil(porcentaje / 100 * len(ordenados)) - 1
    return ordenados[max(0, min(indice, len(ordenados) - 1))]

async def medir_latencia(self, muestras: int = 5):
    await self.get_latest_event()                # calienta la cache
    con = [(await self.get_latest_event()).latency_ms for _ in range(muestras)]
    sin = [(await self.get_latest_event(usar_cache=False)).latency_ms
           for _ in range(muestras)]
    return {"con_cache": _resumen(con), "sin_cache": _resumen(sin)}
\end{lstlisting}

% =====================================================================
\section{Pruebas automatizadas (\arch{tests/clients/})}
% =====================================================================

Se usan \textbf{pytest} + \textbf{pytest-asyncio} + \textbf{respx} (para
simular HTTP sin red). \arch{conftest.py} limpia el estado global
(\arch{\_no\_antes\_de}, cache de modulo) y reemplaza \arch{\_dormir} para
medir las esperas \emph{sin dormir de verdad}.

\begin{table}[H]
\centering
\small
\begin{tabular}{p{7.6cm} p{6.6cm}}
\toprule
\textbf{Prueba} & \textbf{Que demuestra} \\
\midrule
\arch{test\_503\_luego\_200\_reintenta} & Reintenta y termina en OK \\
\arch{test\_429\_respeta\_retry\_after} & Usa \arch{Retry-After}, no el backoff \\
\arch{test\_400\_no\_se\_reintenta} & 400 no se reintenta \\
\arch{test\_404\_lectura\_no\_se\_reintenta} & 404 $\to$ \arch{NoEncontrado} \\
\arch{test\_503\_persistente\_agota\_reintentos} & Agota intentos $\to$ error \\
\arch{test\_timeout\_persistente\_...} & Timeout $\to$ \arch{TimeoutDependencia} \\
\arch{test\_paginacion\_con\_fallo\_...} & Reintenta por pagina \\
\arch{test\_respuesta\_invalida\_listado/lectura} & Pydantic $\to$ \arch{RespuestaInvalida} \\
\arch{test\_degradacion\_parcial\_lecturas} & Un fallo no tumba al resto \\
\arch{test\_salud\_y\_sensor\_existe} & \arch{salud()} y existencia \\
\arch{test\_no\_antes\_de\_compartido} & El bloqueo por IP es global \\
\arch{test\_lecturas\_concurrentes\_...} & Concurrencia respeta el bloqueo \\
\midrule
\arch{test\_encadena\_dos\_llamadas\_y\_usa\_cache} & Feed $\to$ detalle + cache \\
\arch{test\_sin\_detalle\_devuelve\_...} & Evento sin \arch{detail} \\
\arch{test\_feed\_invalido\_...} & Feed invalido $\to$ error \\
\arch{test\_timeout\_lanza\_timeoutdependencia} & Timeout de USGS \\
\arch{test\_obtener\_ultimo\_sismo\_devuelve\_hit} & Funcion de contrato \\
\arch{test\_percentil\_y\_resumen\_latencia} & Matematica del p95 \\
\arch{test\_medir\_latencia\_compara\_...} & Con cache: 1 llamada real \\
\bottomrule
\end{tabular}
\caption{Las 20 pruebas de \arch{tests/clients/}.}
\end{table}

\begin{lstlisting}[caption={Conftest: jitter y bloqueo medidos sin dormir}, language=Python]
@pytest.fixture
def dormidas(monkeypatch):
    registro: list[float] = []
    async def falso_dormir(segundos: float) -> None:
        registro.append(segundos)
    monkeypatch.setattr(us, "_dormir", falso_dormir)
    return registro
\end{lstlisting}

Resultado de correr toda la suite:

\begin{lstlisting}[language=bash, caption={Salida real de pytest}]
$ ./venv/bin/python -m pytest tests/clients -q
....................                                                     [100%]
20 passed in 0.20s
\end{lstlisting}

% =====================================================================
\section{Evidencias reales}
% =====================================================================

Toda la evidencia se genera \textbf{contra el servicio real} (no simulada).

\subsection{Reintentos, \arch{Retry-After} y degradacion parcial}
Fragmento de \arch{evidencias/logs\_clients.txt}
(\arch{--semilla 42 --tasa-fallos 0.4}). Se ve el backoff creciendo, el
\arch{Retry-After} ganando y el 404 que \emph{no} se reintenta.

\begin{lstlisting}[caption={Evidencia real de reintentos (fragmento)}]
intento 1/5, causa=503, espera=0.170s, origen=backoff
intento 2/5, causa=503, espera=0.564s, origen=backoff
intento 3/5, causa=429, espera=2.000s, origen=Retry-After
...
GET /sensores/S-999/lectura "HTTP/1.1 404 Not Found"   <-- NO se reintenta

== obtener_lecturas (degradacion parcial) ==
  lecturas ok: 2  fallidos: 2
  ok:      {'sensor_id': 'S-001', 'tipo': 'temperatura', 'valor': 28.3, ...}
  fallido: {'sensor_id': 'S-003', 'motivo': '...sigue devolviendo 429...'}
  fallido: {'sensor_id': 'S-999', 'motivo': 'No existe el sensor S-999'}

== listar_sensores (recorre toda la paginacion) ==
total sensores: 23
\end{lstlisting}

\subsection{Latencia con y sin cache (C8)}
Corrida real contra USGS (\arch{evidencias/latencia\_cache.txt}).

\begin{table}[H]
\centering
\begin{tabular}{l r r r}
\toprule
\textbf{Modo} & \textbf{Muestras} & \textbf{Media (ms)} & \textbf{p95 (ms)} \\
\midrule
\ok\; Con cache  & 5 & 0.002 & 0.003 \\
\no\; Sin cache  & 5 & 479.529 & 1174.669 \\
\bottomrule
\end{tabular}
\caption{Latencia de la API externa con y sin cache (5 muestras).}
\end{table}

\textbf{Lectura:} la lectura de cache es sub-milisegundo mientras que la
llamada real ronda los 480\,ms de media (p95 $\approx$ 1.17\,s). El TTL
reduce latencia \emph{y} carga sobre la dependencia externa.

% =====================================================================
\section{Hallazgos para el equipo (analisis)}
% =====================================================================

Estos son problemas que detecte al probar el sistema completo y que
corresponden a otras partes:

\begin{table}[H]
\centering
\small
\begin{tabular}{p{2.4cm} p{5.4cm} p{6.4cm}}
\toprule
\textbf{Quien} & \textbf{Hallazgo} & \textbf{Recomendacion} \\
\midrule
Nataly & \arch{Sensor} sin campo \arch{tipo} y con \arch{extra=\quot forbid\quot};
el servicio devuelve \arch{tipo}. & Agregar \arch{tipo: str | None = None} o
relajar a \arch{extra=\quot ignore\quot}. \\
Nataly & \arch{UsgsFeed}/\arch{UsgsProperties} con \arch{extra=\quot forbid\quot}
vs.\ extras reales de USGS. & Cambiar a \arch{extra=\quot ignore\quot} (defensa en
profundidad; yo ya lo tolero). \\
Nataly & Faltan constantes del contrato y \arch{obtener\_o\_calcular} con
\emph{single-flight}. & Alinear \arch{config.py}/\arch{cache.py} al
contrato acordado. \\
Daniela & \arch{NoEncontrado} cae en el 502 en vez de 404. & Agregar
\arch{except NoEncontrado}: \arch{HTTPException(404)}. \\
Daniela & Misma \arch{Idempotency-Key} con cuerpo distinto devuelve la
anterior (deberia ser 422). & Guardar un hash del cuerpo y comparar. \\
Daniela & Lock global en vez de por clave; \arch{/docs} sin
\arch{summary}/\arch{responses}. & Lock por clave y \arch{responses=\{...\}}. \\
Grupal & Faltan \arch{requests.http}, tabla de latencias en README, demo de
idempotencia y deps de test. & Completar entregables. \\
\bottomrule
\end{tabular}
\caption{Hallazgos y recomendaciones (fuera de mi alcance).}
\end{table}

\textbf{Evidencia del 502/504 con la config ``fuerte'':} con
\arch{--tasa-fallos 0.4} y el \arch{max\_retries=3} de \arch{config.py}, las
peticiones a veces agotan intentos antes que los 5 del contrato, y el
endpoint responde 502/504. Es un sintoma de la falta de alineacion de
constantes, no del cliente.

% =====================================================================
\section{Conclusiones}
% =====================================================================

\begin{table}[H]
\centering
\begin{tabular}{l l l}
\toprule
\textbf{Requisito} & \textbf{Estado} & \textbf{Donde} \\
\midrule
A1 Consumo de API externa real       & \ok & \arch{external\_api.py} \\
A2 Validacion Pydantic               & \ok & \arch{external\_api.py} \\
A3 Manejo de errores de la API       & \ok & \arch{external\_api.py} \\
C2 Timeout explicito                 & \ok & \arch{unstable\_service.py} \\
C3 Reintentos selectivos             & \ok & \arch{unstable\_service.py} \\
C4 Backoff + jitter                  & \ok & \arch{unstable\_service.py} \\
C5 \arch{Retry-After} + bloqueo IP   & \ok & \arch{unstable\_service.py} \\
C6 Semaforo de concurrencia          & \ok & \arch{unstable\_service.py} \\
C7 Paginacion resiliente             & \ok & \arch{unstable\_service.py} \\
C8 Latencia media/p95 con y sin cache& \ok & \arch{external\_api.py} \\
Pruebas automatizadas (20)           & \ok & \arch{tests/clients/} \\
Evidencia real                       & \ok & \arch{evidencias/} \\
\bottomrule
\end{tabular}
\caption{Estado final de mi parte.}
\end{table}

Mi parte queda \textbf{completa}: consumo resiliente de las dos
dependencias, cache con metricas de latencia, 20 pruebas en verde y
evidencia real. Lo pendiente esta fuera de mi alcance y se detalla en la
seccion anterior.

% =====================================================================
\appendix
\section{Guion sugerido para la exposicion}
% =====================================================================

\begin{enumerate}
  \item \textbf{Contexto (30\,s):} ``una API agregadora que depende de un
        servicio que falla y de una API publica; mi parte es sobrevivir a
        esos fallos''.
  \item \textbf{Arquitectura (30\,s):} mostrar la figura~\ref{fig:arq}.
  \item \textbf{Excepciones (30\,s):} figura~\ref{fig:exc}; enfatizar la
        compatibilidad con el router.
  \item \textbf{Resiliencia (2--3\,min):} recorrer C2--C7 con un fragmento
        cada uno; mencionar \emph{por que} el timeout de lectura es 2\,s
        (menor que la lentitud simulada).
  \item \textbf{API externa (1\,min):} encadenado feed$\to$detalle, y
        \emph{contar el hallazgo} de \arch{extra=\quot forbid\quot} y como lo resolvi.
  \item \textbf{Latencias (1\,min):} tabla con/sin cache
        ($0.002$\,ms vs $479$\,ms).
  \item \textbf{Pruebas y evidencia (1\,min):} 20 pruebas, log real.
  \item \textbf{Hallazgos (1\,min):} tabla de recomendaciones al equipo.
\end{enumerate}

\section{Comandos de reproduccion}
\begin{lstlisting}[language=bash, caption={Como reproducir todo}]
# 1) Servicio inestable (lo entrega el docente; no vive en el repo)
python servicio_inestable.py --semilla 42 --tasa-fallos 0.4

# 2) Pruebas de los clientes
./venv/bin/python -m pytest tests/clients -q

# 3) Evidencia real de reintentos
PYTHONPATH=. ./venv/bin/python tests/clients/evidencia_real.py

# 4) Medicion de latencia (con y sin cache)
PYTHONPATH=. ./venv/bin/python tests/clients/medir_latencia.py

# 5) Levantar la API y ver /docs
./venv/bin/python -m uvicorn main:app --host 127.0.0.1 --port 8000
\end{lstlisting}

\end{document}
